\begin{abstract}
A hidden attractor, whose basin avoids small neighborhoods of the equilibria, cannot be found from them; when its basin is also small, finding it is a rare-event problem among initial conditions. Measuring cost by the integrations needed to detect the attractor with prescribed confidence, we ask when any method beats uniform sampling. We prove that if everything an algorithm observes, apart from hitting the target, is constant on a set $Q$ containing the target basin, no algorithm can guarantee a gain over uniform sampling larger than about $2.6/\mu(Q)$; adaptivity alone is worth 2.6, escape times 2.6 over the measure of the nonescaping set. For $k$ uncoupled units, this bound and the gain of ideal multilevel splitting along the escape time grow geometrically in $k$ with the same base: rare hidden regimes of systems of units are where adaptive search gains by orders of magnitude. A portfolio of splitting along transient duration, active learning of basin boundaries, and a uniform share $\theta$ provably costs at most about $1/\theta$ times uniform sampling and inherits the cost exponent of each autonomous arm. On coupled copies of a hidden-attractor system, its gain grows from 65 (two copies) to $1.3\times10^4$ (four), with base 14 versus 21 predicted from one copy; on a synthetic family, the cost scales as $p^{-0.45}$ with the basin measure $p$; on bounded systems, where escape times are uninformative, transient settling still gives gains that persist in wall time. Verified enumeration reveals a coupling-born hidden chaotic attractor of two Chua circuits.
\end{abstract}

\begin{keywords}
hidden attractors, multistability, basins of attraction, rare events, multilevel splitting, lower bounds, coupled oscillators
\end{keywords}

\begin{MSCcodes}
37M22, 37C70, 34C28, 65C05, 68W20
\end{MSCcodes}
